\chapter{Polynômes}
\section{Anneaux}
\begin{definition}[Anneau]
Un \footnote{Lecture 2} anneau $R$ est un ensemble $(R,+,\cdot)$ qui vérifie ces conditions
\begin{enumerate}
\item $\forall a,b\in R:a+b=b+a$
\item $\forall a,b,c\in R:a+(b+c)=(a+b)+c$
\item $\exists0\in R:\forall a\in R:a+0=a$
\item $\forall a\in R~\exists(-a)\in R:a+(-a)=0$
\item $\forall a,b,c\in R:a\cdot(b\cdot c)=(a\cdot b)\cdot c$
\item $\exists 1 \in R:\forall a\in R:1\cdot a=a\cdot1=a$
\item $\forall a,b,c\in R:a\cdot(b+c)=a\cdot b+a\cdot c$
\item Si de plus $\forall a,b\in R$ on a $a\cdot b=b\cdot a$ on dit que $R$ est commutatif.
\end{enumerate}
\end{definition}

\begin{theorem}
$$
\forall a\in R:a\cdot0=0
$$
\end{theorem}
\begin{proof}
\begin{align*}
a 0  & =a(0+0) \\
 & =a 0+a 0 \\
0 & =a0 & |+(-a 0)
\end{align*}

\end{proof}
\begin{definition}[Centre de $R$]
$$
Z(R)=\{ r\in R:\forall a\in R:ra=ar \}$$
\end{definition}
\begin{definition}[Diviseur de 0]
Un element de $a\in R\setminus \{ 0 \}$ est un diviseur de 0 s'il existe $b\in R\setminus \{ 0 \}$ tel que $ab=0$ ou $ba=0$.
\end{definition}
\begin{definition}[Anneau intègre]
Un anneau intègre est un anneau commutatif sans diviseur de zero ($1\neq 0$).
\end{definition}
\begin{example}
Montrer que 1 est unique.

\begin{proof}
On pose 1' tel que $1=1'1=1'$.
\end{proof}
\end{example}
\begin{example}
\begin{itemize}
\item $(\mathbf{Z},+,\cdot)$ est un anneau intègre.
\item $(5\mathbf{Z}:=\{ 5z:z\in \mathbf{Z} \},+,\cdot)$ n'est pas un anneau.
\item $(\mathbf{Z}^{2\times2},+,\cdot)$ est un anneau non-commutatif que possède des diviseurs de zero. Pas commutatif car par exemple
$$
\begin{pmatrix}
1 & 0 \\
1 & 1
\end{pmatrix}\begin{pmatrix}
1 & 1 \\
0 & 1
\end{pmatrix}=\begin{pmatrix}
1 & 1 \\
1 & 2
\end{pmatrix}\neq \begin{pmatrix}
1 & 1 \\
0 & 1
\end{pmatrix}\begin{pmatrix}
1 & 0 \\
1 & 1
\end{pmatrix}=\begin{pmatrix}
2 & 1 \\
1 & 1
\end{pmatrix}
$$
et a des diviseurs de zero
$$
\begin{pmatrix}
0 & 0 \\
0 & 1
\end{pmatrix}\begin{pmatrix}
1 & 0  \\
0 & 0
\end{pmatrix}=\begin{pmatrix}
0 & 0 \\
0 & 0
\end{pmatrix}
$$
\end{itemize}
\end{example}
\begin{definition}[Element inversible]
Un element $r\in R$ est inversible s'il existe $(r^{-1})\in R$ tel que $r^{-1}r=1=rr^{-1}$
\end{definition}
\begin{definition}[Ensembles d'unites]
On note l'ensemble des elements inversibles de $R$ par $R^{*}$. $(R^{*},\cdot)$ est un groupe.
\end{definition}

\begin{example}
	Soit $R$ un anneau et on considère $R^{n\times n}$ de centre $Z(R^{n\times n})=\{ aI_{n}:a\in Z(R) \}$. Comme $aI_{n}\in Z(R^{n\times n})$ ca implique que $a\in Z(R)$ car $\forall b\in R$ on a que $aI_{n}\cdot bI_{n}=bI_{n}\cdot aI_{n}$.
\end{example}
\begin{theorem}[Critère de sous-anneau]
Soit $R$ un anneau et $S\subseteq R$, on dit que $S$ est un sous-anneau si $S$ verifie ces conditions:
\begin{enumerate}
\item \label{ssanneau1} $1\in S$
\item \label{ssanneau2} $\forall s,t\in S:s-t\in S$ et $s\cdot t\in S$
\end{enumerate}
\end{theorem}
\begin{theorem}\label{thm22}
Soit $R$ un anneau alors il existe un anneau $R'\supseteq R$ tel que $R$ est un sous-anneau de $R'$ et $1_{R'}=1_{R}$ et ces conditions sont satisfaites:

Il existe $x \in R\setminus R$ tel que
\begin{enumerate} 
\item  $\forall a\in R:ax=xa$
\item Si pour $a_{0},\dots,a_{n}\in R:a_{0}+a_{1}x+\dots+a_{n}x^{n}=0$ alors tous les $a_{i}$ sont egaux a 0.
\end{enumerate}
On appele $x$ une variable ou indeterminee de $R$. 
\end{theorem}
\section{Polynomes}
\begin{definition}[Polynome]
Un polynome sur $R$ est une expression du type:
$$
f(x)=a_{0}+a_{1}x+\dots+a_{n}x^{n}
$$
ou $a_{i}\in R,i=0,1,\dots,n$ et $x \in R'\setminus R$ est l'indeterminee.

On appelle $R[x]\subseteq R'$ l'ensemble des polynomes sur $R$.
\end{definition}
\begin{theorem}
$R[x]$ est un sous-anneau de $R'$. Si $R$ est commutatif alors $R[x]$ l'est aussi et si $R$ ne possede des diviseurs de zero alors $R[x]$ non plus.
\end{theorem}
\begin{proof}
On a que $1\in R[x]$ donc on verifie \ref{ssanneau1}. $\checkmark$ Donc il nous reste a montrer que $\forall f(x),g(x)\in R[x]$: 
\begin{itemize}
\item $f(x)-g(x)\in R[x]$
\item $f(x)\cdot g(x)\in R[x]$
\end{itemize}
pour verifier \ref{ssanneau2}.
\begin{itemize}
\item Soient $f(x)=a_{0}+a_{1}x+\dots+a_{n}x^{n}$ et $g(x)=b_{0}+b_{1}x+\dots+b_{n}x^{n}$ alors
$$
f(x)-g(x)=(a_{0}-b_{0})+(a_{1}-b_{1})x+\dots+(a_{n}-b_{n})x^{n}
$$
donc $f(x)-g(x)\in R[x]$ car $(a_{i}-b_{i})\in R$.
\item

\begin{align}
f(x)g(x) & =(a_{0}+a_{1}x+\dots +a_{n}x^{n})(b_{0}+b_{1}x+\dots+b_{n}x^{n}) \\
 & =a_{0}b_{0}+a_{1}b_{0}x +a_{2}b_{0}x^{2}+\dots \\
 & =\sum_{i=0}^{2n} \left( \sum_{k+j=i} a_{k}b_{j}\right)x^{i} \in R[x] \label{comm1}
\end{align}
\end{itemize}
La commutatitivite de $R$ implique celle de $R[x]$ grace a la formule \ref{comm1}. Pour montrer que $R[x]$ ne possede des diviseurs de zero si $R$ non plus on pose deux polynomes non-nuls $p(x)=a_{0}+\dots+a_{m}x^{m}$ et $g(x)=b_{0}+\dots+b_{\mu}x^{\mu}$ avec $a_{m}\neq0$ et $b_{\mu}\neq0$. Alors leur produit par \ref{comm1} aura au moins un terme non-nul $a_{m}b_{\mu}x^{\mu+m}$.
\end{proof}
\begin{example}
\begin{itemize}
\item $p(x)=3+x^{2}+5x^{4}\in \mathbf{Z}[x]$
\item $f(x)=\begin{pmatrix} 1 & 0 \\  1 & 1\end{pmatrix}+\begin{pmatrix}3 & 3 \\ 2 & 1\end{pmatrix}x^{3}\in \mathbf{Z}^{2\times2}[x]$
\end{itemize}
\end{example}

\begin{proposition}
Soient deux polynomes $f(x)=a_{0}+a_{1}x+\dots+a_{n}x^{n}$ et $g(x)=b_{0}+\dots+b_{m}x^{m}$ dans $R[x]$. On dit que $f,g$ sont egaux si et seulement si $a_{0}=b_{0},a_{1}=b_{1},\dots$
\end{proposition}
\begin{proof}
\begin{itemize}
\item La direction "$\Leftarrow$" est claire
\item Pour la direction "$\Rightarrow$": Sans perte de generalite on pose $n=m$. Si $f(x)=g(x)$ alors 

\begin{align*}
 a_{0}+a_{1}x+\dots+a_{n}x^{n} & =b_{0}+b_{1}x+\dots+b_{n}x^{n} \\
\Leftrightarrow(a_{0}-b_{0})+(a_{1}-b_{1})x+\dots.+(a_{n}-b_{n})x^{n} & =0 \\
\Leftrightarrow(a_{i}-b_{i}) & =0 \\
\Leftrightarrow a_{i} & =b_{i}
\end{align*}

\end{itemize}
\end{proof}
\begin{definition}
Soit $f(x)=a_{0}+a_{1}x+\dots+a_{n}x^{n}\in R[x]$ alors $a_{i}\in R$ est le $i$-eme coefficient de $f(x)$.
\end{definition}
\begin{definition}[Degre d'un polynome]
Soit $f(x)=a_{0}+\dots+a_{n}x^{n}\in R[x]$ un polynome. On note de le degre de $f$:
$$
\deg f =\max\{i\in \mathbf{N}:a_{i}\neq 0 \}
$$
\end{definition}
\begin{example}
Le polynome $f(x)=x^{2}+2\in \mathbf{Q}[x]$ est $\deg f=2$.
\end{example}
\begin{definition}[Coefficient dominant]
On dit que $a_{\deg f}$ est le coefficient dominant de $f(x)$.
\end{definition}
\begin{remark}
Si $f(x)=0$ on dit que $\deg f=-\infty$.
\end{remark}
\begin{theorem}
Soient $f,g\in R[x]\setminus \{ 0 \}$ tels que le coefficient dominants de $f$ ou $g$ n'est pas un diviseur de zero alors:
$$
\deg (fg)=\deg f+\deg g
$$
\end{theorem}
\begin{proof}
Soit $f(x)=a_{0}+\dots+a_{n}x^{n}$ et $g(x)=b_{0}+\dots+b_{m}x^{m}$ ou $a_{n}\neq 0$ et $b_{m}\neq 0$. Alors $a_{n}b_{m}\neq 0$ et par \ref{comm1} on a que $\deg(f+g)=n+m=\deg f +\deg g$.
\end{proof}

\begin{definition}[Evaluation]
Soit $p(x)=a_{0}+\dots+a_{n}x^{n}\in R[x]$ un polynome alors $p(x)$ induit une application

\begin{align*}
f_{p}:R & \to R \\
r & \mapsto p(r)
\end{align*}

qui associe $p(r)=a_{0}+\dots+a_{n}r^{n}$ a $r$. 
\end{definition}
\begin{remark}
Deux polynômes différentes peuvent donner les memes applications. Mais si $R=K$ est un corps infini alors $f_{p}\neq f_{g}$ pour tous $f,g\in K[x]$.
\end{remark}

